\section{Data and Sample}
\label{sec:data}

\subsection{Sources}

Three administrative datasets, January \BPsampleStartYear{} through
December \BPsampleEndYear. The first is the universe of bid-level
records from the BEC electronic procurement platform for SES/SP
purchases of medical supplies (BEC Group~\BPbecGroup). Each
observation is a purchase-offer-item (POI), recording the reference
price set during the planning phase, the winning bid price, the
number of bidding firms, the quantity ordered, the modality
(\textit{preg\~ao eletr\^onico} or \textit{convite}), and the
tender outcome.

The second is a purchase-type classifier that assigns each POI to
one of three regimes: ordinary, administrative-request, or litigated.
Brazilian law requires the tender notice for any urgent purchase to
reference the underlying court order or administrative-channel
request, and the classifier exploits this regulatory disclosure
through a regular-expression algorithm applied to free-text notices.
We validate it on a stratified hand-labeled sample of
\BPregexNotices{} notices ($\sim$\BPregexPerClass{} per predicted
class); detailed metrics are in the Online
Appendix.\footnote{The regex relies on literal court-order
references whose presence is independent of the price outcomes we
estimate, so any residual misclassification attenuates rather than
biases our coefficients.}

The third is the SES/SP S-CODES database of health-related
litigation, which provides SUS-formulary status and injunction
enforcement data for each litigated request.

The binary outcome \emph{tender success} equals one if the POI
closes with an accepted winning bid within its procurement order and
zero otherwise; re-tendered items enter as new POIs and are not
counted retroactively as successes. All continuous variables are
winsorized at the \BPwinsorLow{}st and \BPwinsorHigh{}th
percentiles; robustness to alternative winsorization is in the
Appendix.

\subsection{Sample and descriptive patterns}

The raw data comprise \BPnPOIfull{} POI observations across
\BPnItemsFull{} distinct items, \BPnOrdersFull{} purchase orders,
and \BPnPBUsObserved{} PBUs (Online Appendix Table~A.1). Our
analysis sample restricts to items for which at least one ordinary
and at least one litigated purchase is observed, ensuring within-item
comparability; this yields \BPnAnalysisSample{} observations
covering \BPnItemsAnalysis{} items.\footnote{Sample sizes vary
slightly across tables due to missing values in specific dependent
variables and the restriction to winners for price regressions.}
Price regressions further restrict to winning bids (\BPnWinners{}
observations).

Table~\ref{tab:descriptive_balance} reports sample means by purchase
type on winning bids. Three patterns anticipate the empirical
results. \emph{Higher litigated prices.} Log negotiated prices are
higher in litigated than in ordinary purchases (\BPmeanLogNegLit{}
vs.\ \BPmeanLogNegOrd; item fixed effects absorb the level mix).
\emph{Larger administrative orders.} Administrative orders are
markedly larger than litigated ones (\BPmeanQtyAdm{} vs.\
\BPmeanQtyLit{} units on average), the bulk-discount disadvantage
that drives the mechanical channel in Section~\ref{sec:results}.
\emph{Thinner urgent competition with higher success.} The number of
bidding firms is lower in urgent than in ordinary tenders
(\BPmeanFirmsLit{} for litigated vs.\ \BPmeanFirmsOrd{} for
ordinary), yet tender success rates are higher in urgent regimes
(\BPsuccessLit{} litigated vs.\ \BPsuccessOrd{} ordinary)---officials
accept worse terms to ensure compliance.

\begin{table}[ht]
\centering
\caption{Sample means by purchase type.}
\label{tab:descriptive_balance}
\begin{threeparttable}
\begin{tabular}{lccc}
\toprule
 & Ordinary & Administrative & Litigated \\
\midrule
Observations (winners) & 343{,}393 & 16{,}985 & 46{,}212 \\
Log reference price & 2.109 & 0.960 & 1.578 \\
Log negotiated price & 1.596 & 0.543 & 1.050 \\
Quantity (units) & 56{,}280 & 794{,}534 & 61{,}903 \\
Number of bidding firms & 3.159 & 2.394 & 2.191 \\
Tender success rate & 84.1\% & 85.7\% & 90.7\% \\
\bottomrule
\end{tabular}
\begin{tablenotes}[flushleft]\footnotesize
\item Means computed on winners (POIs with an accepted winning bid). Reference and negotiated prices are in logs; quantity in units; tender success rate is the share of POIs that close with an accepted winning bid. Sample period: 2009--2019. Source: BEC-SP, SES/SP S-CODES.
\end{tablenotes}
\end{threeparttable}
\end{table}



\subsection{The under-the-gun subsample}

For the under-the-gun analysis, we construct a subsample of urgent
purchases (administrative and litigated) for items with both regimes
present, yielding \BPnUTG{} observations (Online Appendix
Table~A.2). The tightest specifications identify off
item$\times$year-month cells in which both regimes coexist: of
\BPutgCellsItemYM{} such cells, \BPutgCellsBoth{}
(\BPutgCellsBothPct) carry within-cell variation. The remaining
cells are absorbed as singletons under the fixed-effects design.
Both-types cells are not a representative subsample---they tilt
toward higher-volume, more-reordered medications (Online Appendix
Table~\ref{tab:both_types_cells}); the bounded estimate therefore
speaks to the more liquid items in the urgent panel, a restriction
we revisit in the welfare bound. Online Appendix complements the
balance with geographic evidence on the spatial distribution of
administrative demand, which mirrors the litigation patterns of
Section~\ref{sec:institutional}.
